\chapter{Derandomisation}

\section{Introduction}

Randomised algorithms often achieve remarkable simplicity and efficiency. However, they raise a natural 
question $-$ how much randomness is really necessary? Recall that we stated that for 
our purposes, pseudo-randomness is true randomness. In this chapter, we will finally differentiate between 
the two! Since true randomness is often expensive or unavailable, it is natural to ask whether one can 
reduce, or even eliminate, the amount of randomness required. The study of such questions forms the 
subject of \emph{derandomisation}. 

\section{Independence}

\subsection{Mutual Independence}
A set of random variables $S=\{X_1,X_2,\dots,X_n\}$ are said to be \emph{mutually independent} if for 
every subset $s\subseteq S$, we can write the probability of a particular assignment of values as the product 
of the individual probabilities. 

\[
\{X_{i_1},X_{i_2},\dots,X_{i_t}\}\subseteq\{X_1,X_2,\dots,X_n\}
\]
\[
\implies 
\Pr[X_{i_1}=a_1,\dots,X_{i_t}=a_t]
=
\prod_{j=1}^{t}\Pr[X_{i_j}=a_j]
\]

This definition is the strongest possible definition of independence and is the default meaning of the 
term ``independent''.

\subsection{k-Wise Independence}
A set of random variables $X_1,X_2,\dots,X_n$ are said to be \emph{$k$-wise independent} if every subset of 
at most $k$ variables is mutually independent.\\


\noindent By definition it follows that $k+1$-wise independence implies $k$-wise independence. Of course, 
$n$-wise independence is the same as mutual independence of the whole set and is thus the strongest 
version of $k$-wise independence. The case $k=2$ is the weakest version of $k$-wise independence and 
is called \emph{pairwise independence}.

\section{Generating Pairwise Independent Bits}

Let us say that we wish to generate a pairwise independent set of $n=2^k$ random bits. Our source of randomness 
is a set of $m$ truly random bits $R_1,R_2,\dots,R_m$, where by truly random what we really 
mean is mutually independent. Our goal will be to achieve the same $n$ with minimum usage of truly 
random bits since they are expensive to generate.

\subsection{Linear Map}

Let $n \in \mathbb K=\{0,1,\dots,2^k-1\}$. Consider the finite field $\mathbb F_{2^k}$ $-$ recall that in the 
standard construction of finite fields, the elements are polynomials and the operations are taken modulo a 
particular polynomial called the irreducible polynomial. We assign indices to the elements of $\mathbb F_{2^k}$ as 
$0,1,\dots,2^k-1$ (in any order) for our convenience. Hence the elements of $\mathbb F_{2^k}$ are 
$\{p_0,p_1,\dots,p_{2^k-1}\}$.\\ 

\noindent Choose uniformly random values $a,b \in \mathbb F_{2^k}$ and define the function 
$T : \mathbb K \to \{0,1\}$ as $T(i) = \operatorname{const}(a\cdot p_i+ b)$. The function $T$ basically 
uses a linear mapping to convert the element $p_i$ to a random element and then extracts the 
constant term of the resulting polynomial. It is quite straightforward to prove that $a\cdot p_i+b$ for a 
given $i$ is uniformly distributed over $\mathbb F_{2^k}$. Consequently, $T(i)$ is uniformly distributed over 
$\{0,1\}$.

\[
\implies \Pr(T(i)=0)=\Pr(T(i)=1)=\frac12
\]

\noindent Now consider the random bits generated by $T$ as the set of random variables 
$\{T(0),T(1),\dots,T(2^k-1)\}$. With some thought, it can be proved that this set is pairwise independent.

\subsection*{Cost}
To generate independent and uniformly distributed $a$ and $b$, we need $2k$ truly random bits. 
This is because each of $a$ and $b$ can take $2^k$ values and we can identify a particular value in the 
finite field $\mathbb F_{2^k}$ using the binary representation of the index of the element. In terms of $n=2^k$, 
we needed $2\log_2 n$ truly random bits to generate $n$ pairwise independent random bits. This 
is already quite efficient in terms of the number of truly random bits used, but can we do better? 

\subsection{XOR Construction}
Choose $m$ truly random bits $R_1,R_2,\dots,R_m$. For every 
non-empty subset $S\subseteq\{1,2,\dots,m\}$, define the random bit

\[
X_S=\bigoplus_{i\in S}R_i,
\]

where $\oplus$ denotes the XOR operation. It is fairly straightforward to prove that $X_S$ is equally 
likely to be $0$ or $1$.

\[
\Pr(X_S=0)=\Pr(X_S=1)=\frac{1}{2}
\]

To prove pairwise independence, consider two distinct subsets
$S$ and $T$. Since $S\neq T$, there exists a bit $R_z$ belonging to exactly
one of them. Without loss of generality, let $z\in S - T$.
Fix all random bits except $R_z$. Flipping $R_z$ changes $X_S$ but leaves $X_T$ unchanged. So for every 
fixed value of $X_T$, the bit $X_S$ is equally likely to be $0$ or $1$.

\[
\Pr(X_S=v \;| \; X_T=w)=\frac12=\Pr(X_S=v)
\]

\[
\implies 
\frac{\Pr((X_S=v) \cap (X_T=w))}{\Pr(X_T=w)}=\Pr(X_S=v)
\]

\[
\implies 
\Pr((X_S=v) \cap (X_T=w))=\Pr(X_S=v)\cdot\Pr(X_T=w)
\]

\subsection*{Cost}
Since we generated $n=2^m-1$ pairwise independent random bits using $m$ truly random bits, the cost for a fixed 
$n$ is $m=\log_2(n+1)$. This again begs the question, can we do better? Turns out, we really can not!

\subsection{Minimum Cost}

Let us assume the existence of a set of $n$ pairwise independent random bits 
generated using $m$ truly random bits. Without loss of generality, consider the generator to be a function
$F:\{0,1\}^m \to \{0,1\}^n$ from every instance of the $m$ truly random bits to a particular instance of 
the $n$ pairwise independent random bits. For our convenience, we represent an instance of the $n$ pairwise 
independent random bits as a row matrix of dimension $1\times n$.\\

Since there are $2^m$ possible instances of the $m$ truly random bits, we can consider the matrix $M$ of 
dimension $2^m \times n$ where the $i$-th row of $M$ is the output of $F$ for the $i$-th instance 
of the $m$ truly random bits, indexed arbitrarily (say, lexicographically). For $n=4$ and $m=3$ and 
representing the $i$-th instance of the $m$ truly random bits as $R_i$, the matrix $M$ is given by

\[
M=
\begin{bmatrix}
F(R_0=(0,0,0))\\
F(R_1=(0,0,1))\\
F(R_2=(0,1,0))\\
F(R_3=(0,1,1))\\
F(R_4=(1,0,0))\\
F(R_5=(1,0,1))\\
F(R_6=(1,1,0))\\
F(R_7=(1,1,1))\\
\end{bmatrix}
=
\begin{bmatrix}
0&0&0&0\\
0&0&1&1\\
0&1&0&1\\
0&1&1&0\\
1&0&0&1\\
1&0&1&0\\
1&1&0&0\\
1&1&1&1
\end{bmatrix}
\]

\noindent Interestingly, the function $F$ in the above example is indeed a valid generator for $n=4$ pairwise independent 
random bits using $m=3$ truly random bits. This can be verified by checking that the pairs 
$(0,0),(0,1),(1,0)$ and $(1,1)$ occur equal number of times if we consider any two columns side-by-side since 
it is equivalent to verifying that the probability of a particular pair of bits is $1/4$ which is the exact 
probability required for pairwise independence. Of course, one may question the need for a matrix 
representation of this? The reason why we formulated this in terms of matrices is because we 
will exploit the properties of the matrix operator \textbf{rank}.\\

Consider $M$ for some $m$ and $n$. Replace every entry of $M$ which is $0$ with $-1$. From our earlier discussion, 
we know that this means that for every pair of columns side-by-side, the number of occurences of $(1,1)$, 
$(1,-1)$, $(-1,1)$, and $(-1,-1)$ will be equal. 

\[
M'
=
\begin{bmatrix}
-1&-1&-1&-1\\
-1&-1&1&1\\
-1&1&-1&1\\
-1&1&1&-1\\
1&-1&-1&1\\
1&-1&1&-1\\
1&1&-1&-1\\
1&1&1&1
\end{bmatrix}
\]

So if we take the dot product of those two columns considering them as vectors, the resulting value 
will be $0$.

\[
C_i \cdot C_j = \sum_{k=1}^n C_{ik} \cdot C_{jk} = 2^{m-2} \cdot (1 \times 1) + 2^{m-2} \cdot (1 \times -1) + 2^{m-2} \cdot (-1 \times 1) + 2^{m-2} \cdot (-1 \times -1) = 0
\]

So any two columns of $M$ are orthogonal, implying that the rank of $M$ is the number of columns of $M$ which is 
$n$. Since the rank of a matrix is at most the minimum of the number of rows and columns, 

\[
\operatorname{rank}(M) \leq \min(n,2^m) \leq 2^m
\]

\[
\operatorname{rank}(M)=n
\implies n \leq 2^m 
\implies m \geq \log_2(n)
\]

Is this the best bound? With our clever XOR construction, we had achieved something slightly worse than this. 
It turns out that with a slight modification to the proof, we can achieve the stricter bound whose 
equality occurs at the XOR construction. Consider prepending a column of all ones to the matrix $M'$. 

\[
M''
=
\begin{bmatrix}
1&-1&-1&-1&-1\\
1&-1&-1&1&1\\
1&-1&1&-1&1\\
1&-1&1&1&-1\\
1&1&-1&-1&1\\
1&1&-1&1&-1\\
1&1&1&-1&-1\\
1&1&1&1&1
\end{bmatrix}
\]

The condition that any two columns of $M'$ are orthogonal still holds trivially because each of the original 
columns have equal number of $1$ and $-1$ entries. In this version of the proof, 
we conclude that the rank of $M'$ is $n+1$.

\[
\operatorname{rank}(M) \leq \min(n+1,2^m) \leq 2^m
\]

\[
\operatorname{rank}(M)=n+1
\implies n+1 \leq 2^m 
\implies m \geq \log_2(n+1)
\]


\section{Bipartite Subgraph Yet Again}

Recall the probabilistic algorithm of repeatedly generating a random 2-coloring of the
vertices and checking the number of bichromatic edges. The expected number of bichromatic
edges was shown to be at least $m/2$, implying the existence of a coloring with at least
$m/2$ bichromatic edges. However, if we carefully examine the proof, we notice that we
never actually used the mutual independence of the random bits. For every edge $(u,v)$,
the only property we required was that the colors assigned to $u$ and $v$ were pairwise
independent. This ensured that

\[
\Pr(\text{$u$ and $v$ receive different colors})=\frac12,
\]

which in turn implied that the expected number of bichromatic edges is at least $m/2$. 
Therefore, instead of assigning colors using $|V|$ mutually independent random
bits, we may use our pairwise independent bit generator. From the previous section, we
know that if we consider the smallest $k$ such that $2^k \geq |V|+1$, this 
requires only $k$ truly random bits. Explicitly, $k$ is equal to $\lceil\log_2(|V|+1)\rceil$. 
So instead of proving the existence of a valid coloring in the complete set of colorings of size $2^{|V|}$, 
we proved the existence of a valid coloring in the set of colorings of 
size $2^{\lceil\log_2(|V|+1)\rceil} \le 2(|V|+1)$. This is a much stronger statement than the original one!\\


\noindent This observation leads us to a completely derandomised and thus completely deterministic 
algorithm. Since our search space is polynomial in size, we can simply enumerate all possible colorings and 
we are guaranteed to find a coloring satisfying our requirement of at least $m/2$ bichromatic edges. 

\subsection{The Algorithm}

\begin{center}
\begin{minipage}{0.9\linewidth}
\begin{framed}
\begin{Verbatim}

n=|V|;
k=ceil(log(n+1.base=2));
vector bestColoring(n);
bestSize=0;
vector all_seeds=GenerateAllSeeds(k);

for (seed in all_seeds)
{
    coloring=GenerateColoring(s)
    cutSize=CountBichromaticEdges(coloring)

    if (cutSize>bestSize)
    {
        bestSize=cutSize
        bestColoring=coloring
    }
}
print(bestColoring);
\end{Verbatim}
\vspace{-4mm}
\end{framed}
\end{minipage}
\end{center}

\subsection{Time Complexity}
\texttt{GenerateAllSeeds} takes $\Theta(n=|V|)$ time. In each iteration, \texttt{GenerateColoring} 
takes $\Theta(n)$ time and \texttt{CountBichromaticEdges} takes $\Theta(m=|E|)$ time. Since there are $n$ such iterations, 
overall time complexity is $\Theta(n(n+m))$.


\subsection{Method Of Conditional Expectation}
Another way to derandomise the algorithm is to use the method of conditional expectation. 
Instead of making random choices, we assign vertices one by one smartly, trying to maximise the 
number of bichromatic edges. Suppose some vertices have already been assigned and the remaining vertices are
still unassigned. Let X be the random variable denoting the number of bichromatic edges.
Let L and R be the bipartitions.\\


Denote the expected value of X given the first \(i\) assignments by:
\[
\mathbb E[X \mid i]
\]


For the $i+1$-th assignment of vertex \(v\), we have two choices:
\begin{itemize}
\item place \(v\) in \(L\),
\item place \(v\) in \(R\).
\end{itemize}

The current conditional expectation is the average of the expectations obtained
from these two choices. Therefore, at least one choice must yield a conditional
expectation that is at least as large as the current one.

\[
\max \left( \mathbb E[X \mid \text{i, v in L}], \mathbb E[X \mid \text{i, v in R}] \right) >= \frac{1}{2} \left(\mathbb E[X \mid \text{i, v in L}] + \mathbb E[X \mid \text{i, v in R}] \right) = \mathbb E[X \mid i]
\]


We simply assign the next vertex such that it would add the most 
bichromatic edges in this step -- that choice yields a better expectation for the future.
\vspace{0.5em}

\[
\text{Since I have ensured that }\mathbb E[X \mid i+1] \ge \mathbb E[X \mid i]
\]

\[
\implies \mathbb E[X \mid \text{n assignments}] \ge \mathbb E[X \mid \text{0 assignments}]
\]

\[
\implies X_{final} \ge \frac{m}{2}
\]

\vspace{0.5em}

\noindent\textbf{Why did that work?}\\


\noindent Initially $\mathbb E[X]=\frac{m}{2}$. At every step we choose an 
assignment that does not decrease the conditional
expectation\footnote{This is by no means a proof for optimality -- it does not 
maximise the number of bichromatic edges.}. Hence the conditional expectation never drops below \(m/2\)!\\


\noindent After all vertices have been assigned, no randomness remains. The conditional
expectation is then equal\footnote{If you are feeling tricked somehow, 
I relate fully! But it is true, we designed this deterministic algorithm and simultaneously proved that 
it produces at least \(m/2\) crossing edges using conditional expectation!} to the 
actual number of crossing edges produced by the algorithm.

\[
\implies
X_{final} \ge \frac{m}{2}.
\]

\noindent Thus we obtain a deterministic algorithm that always constructs a bipartite
subgraph containing at least half of the edges.\\


\noindent\textbf{Why does this not work for all problems?}\\


\noindent The properties of this particular problem we needed for this to work were -- 

\noindent \begin{enumerate}
\item $E[X \mid 1] = m/2$, regardless of which assignment that was -- symmetry between the left and right sides.
\item It was obvious which assignment would maximise $E[X \mid i+1]$ given $E[X \mid i]$.
\end{enumerate}

\noindent Both these properties do not hold, for example, in the Dominating Set problem. Depending 
on your first choice (to put in the dominating set or dominated set), you might get a 
very good dominating set in expectation or a very bad one. An extreme counter example is a star shaped\footnote{A star shaped graph is a graph with 
one central node and \(n-1\) nodes connected to it and only it.} 
graph, where if you put the central node in dominated set you will get a dominating set of size \(n-1\), 
but if you put it in dominating set you will get a dominating set of size 1. 
The second property also fails because in further choices it is not 
clear how to choose which set to put the node in to maximise the conditional expectation.




\clearpage
\section{Problems}
\begin{enumerate}
\item[\textbf{1}]
Consider the polynomial map $P : \mathbb F_{2^n} \to \mathbb F_{2^n}$ defined by 
$P(i)=a_{k-1}\cdot p_i^{k-1} + a_{k-2}\cdot p_i^{k-2} + \cdots + a_1\cdot p_i + a_0$. For a random initialisation 
of $a_0,a_1,\dots,a_{k-1}$, consider the set of random variables $\mathbb G= \{P(0),P(1),\dots,P(2^{n}-1)\}$.

\begin{enumerate}[label=\textbf{(\alph*)}]
\item Show that for any $i$, the random variable $T(i)$ is uniformly distributed over $\mathbb F_{2^n}$ 

\item Construct a system of equations and show existence and uniqueness of the $k$-tuple of coefficients 
$a_0,a_1,\dots,a_{k-1}$ given a particular assignment of values to any $k$-sized subset of 
$\mathbb G$. Show that $\mathbb G$ is $k$-wise independent by showing that the probability 
of a particular assignment of values to any $k$-sized subset of $\mathbb G$ is thus $\frac{1}{2^k}$.

\[
\Pr \left(P(i_1)=v_1,\dots,P(i_k)=v_k \right)
=
\frac{1}{2^k}
=
\prod_{j=1}^{k} \Pr \left(T(i_j)=v_j\right)
\]

\item Show that $\mathbb G$ is not $(k+1)$-wise independent by considering an arbitrary subset of $\mathbb G$ 
of size $k+1$ and constructing a non-trivial formula for the $(k+1)$-th element in terms of the other 
$k$ elements, showing that the probability of a particular assignment of values to the $(k+1)$-th element 
can be $0$ in case the formula is not satisified, disproving $(k+1)$-wise independence.

\end{enumerate}

\item[\textbf{2}]
Recall that a set of $n$ pairwise independent bits generated from $m$ 
truly random bits satisfies $m \geq \log_2(n+1)$, obtained by showing that the columns of 
the associated $\pm 1$ matrix $M'$ (together with a prepended all-ones column) are pairwise 
orthogonal and hence linearly independent. In this problem, we wish to establish a general bound 
for $k$-wise independence. Let $F$, $M$, $M'$, $M''$ be defined as in the text.\\

\noindent Let $\mathcal S = \{S \subseteq \{1,\dots,n\} : |S|\leq k/2\}$ be set of all sets containing at most 
$k/2$ (including $S=\emptyset$). Define the \emph{XORed matrix} $N$ to be the 
$2^m \times |\mathcal S|$ matrix with one column for each 
$S \in \mathcal S$ constructed by taking XOR of the columns element-wise. The XOR of an empty 
set is taken as $1$ for our convenience. Formally, 

\[
N_{r,S} = \bigoplus_{i \in S} M_{r,i}, \;\; N_{r,\phi}=1
\]

Let $N'$ be the corresponding $\pm 1$-valued matrix obtained from $N$ by replacing $0$'s with $-1$'s and define 
$C_S$ be the column generated by the set $S$.

\begin{enumerate}[label=\textbf{(\alph*)}]

\item Show that the number of columns in $N'$ is $|\mathcal S|= \sum\limits_{i=0}^{\lfloor k/2 \rfloor} \binom{n}{i}$

\item For any two distinct subsets $S,T \in \mathcal S$, both of size at most $k/2$, show that the 
dot product of the columns $C_S$ and $C_T$ of $M'$ is zero.

\[
C_S \cdot C_T = 0
\]

\emph{(Hint: consider the symmetric difference $S \triangle T$)}

\item Using the result above, show that

\[
\operatorname{rank}(N') = |\mathcal S| = \sum_{i=0}^{\lfloor k/2 \rfloor} \binom{n}{i}
\]

\item Use the result above and an inequality on the rank of a matrix to show that

\[
m \;\geq\; \log_2 \left( \sum_{i=0}^{\lfloor k/2 \rfloor} \binom{n}{i} \right)
\]

\item Show that for $k=2$, N=M and N'=M'' by noticing that our convention of taking empty set XOR as 1 
automatically created a column of all 1's in N'. Additionally, verify that our new result matches our original 
conclusion for $k=2$.

\[
k=2 \implies m \ge \log_2(n+1)
\]

\end{enumerate}
\end{enumerate}